\section{Architecture \& Implementation}
The \texttt{QRTTurbulenceOptimizer} is implemented as an institutional-grade PyTorch \texttt{Optimizer} subclass within \texttt{nrc\_ai/qrt\_optimizer.py}.

\subsection{Production PyTorch Implementation}
\begin{lstlisting}[language=Python]
from typing import Callable, Iterable, Optional, overload
import torch
from torch.optim import Optimizer
from nrc.math import qrt_damping

class QRTTurbulenceOptimizer(Optimizer):
    def __init__(
        self,
        params: Iterable[torch.Tensor],
        lr: float = 1e-3,
        beta1: float = 0.9,
        beta2: float = 0.999,
    ):
        defaults = {"lr": lr, "beta1": beta1, "beta2": beta2}
        super(QRTTurbulenceOptimizer, self).__init__(params, defaults)

    @torch.no_grad()
    def step(self, closure: Optional[Callable[[], float]] = None) -> Optional[float]:
        loss = None
        if closure is not None:
            with torch.enable_grad():
                loss = closure()

        for group in self.param_groups:
            lr, beta1, beta2 = group["lr"], group["beta1"], group["beta2"]

            for p in group["params"]:
                if p.grad is None:
                    continue

                grad = p.grad
                state = self.state[p]

                if len(state) == 0:
                    state["step"] = 0
                    state["momentum"] = torch.zeros_like(p)
                    state["turbulence"] = torch.zeros_like(p)

                state["step"] += 1
                m, v = state["momentum"], state["turbulence"]

                # Exponential momentum and turbulence accumulation
                m.mul_(beta1).add_(grad, alpha=1.0 - beta1)
                v.mul_(beta2).addcmul_(grad, grad, value=1.0 - beta2)

                bias_correction1 = 1.0 - beta1 ** state["step"]
                bias_correction2 = 1.0 - beta2 ** state["step"]

                m_hat = m / bias_correction1
                v_hat = v / bias_correction2

                # QRT Continuous Damping Engine
                sqrt_turbulence = torch.sqrt(v_hat)
                sqrt_turbulence_np = sqrt_turbulence.detach().cpu().numpy()
                qrt_damped_friction_np = qrt_damping(sqrt_turbulence_np)
                qrt_damped_friction = torch.as_tensor(
                    qrt_damped_friction_np, device=v_hat.device, dtype=v_hat.dtype
                )

                update_step = (m_hat / (qrt_damped_friction + 1e-8)) * lr
                p.sub_(update_step)

        return loss
\end{lstlisting}

\subsection{Pipeline Integration}
\texttt{QRTTurbulenceOptimizer} plugs directly into standard PyTorch training loops, substituting for \texttt{torch.optim.AdamW} with zero modification to loss functions or model definitions.
